\documentclass[10pt,a4paper]{amsart}
\usepackage[margin=19mm]{geometry}
\usepackage{amsmath,amssymb,mathtools}
\usepackage[scr=rsfs,cal=euler]{mathalfa}
\usepackage{xfrac}
\usepackage[unicode,hidelinks,bookmarks=false]{hyperref}
\newcommand{\ie}{i.e.\ }
\newcommand{\cf}{cf.\ }
\newcommand{\resp}{respectively\ }
\newcommand{\Subsection}{\subsection}
\providecommand{\nobreakdash}{\nobreak\hbox{-}}
\setlength{\emergencystretch}{2em}
\multlinegap0pt
\usepackage{bm}
\usepackage{bbm}
\usepackage{dsfont}
\usepackage{xspace}
\usepackage{cancel}
\usepackage{mathtools}
\usepackage{cleveref}
\usepackage{amsthm}
\makeatletter
\@ifundefined{defi}{\newtheorem{defi}{Defi}}{}
\@ifundefined{claim}{\newtheorem{claim}[defi]{Claim}}{}
\@ifundefined{claimproof}{\newtheorem{claimproof}{Claimproof}}{}
\@ifundefined{cor}{\newtheorem{cor}[defi]{Corollary}}{}
\@ifundefined{lem}{\newtheorem{lem}[defi]{Lemma}}{}
\@ifundefined{thm}{\newtheorem{thm}[defi]{Theorem}}{}
\makeatother
\makeatletter
\newcommand*{\abs}[1]{\lvert #1\rvert}
\expandafter\ifx\csname claimproof\endcsname\relax
\newenvironment{claimproof}[1][\myproofname]{\begin{proof}[#1]\renewcommand*{\qedsymbol}{\(\diamondsuit\)}}{\end{proof}}
\fi
\expandafter\ifx\csname customthm\endcsname\relax
\newenvironment{customthm}[1]
 {\renewcommand\theinnercustomthm{#1}\innercustomthm}
 {\endinnercustomthm}
\fi
\newcommand{\diam}{diam}
\makeatother
\title[Uniform \v{S}olt\'es' hypergraphs and \v{S}olt\'es' weighted]{Uniform \v{S}olt\'es' hypergraphs and \v{S}olt\'es' weighted graphs}
\author{Stijn Cambie, Ajay Tiwari}
\date{Source: arXiv:2506.07511v1 (CC BY 4.0)}
\begin{document}
\maketitle

\noindent\emph{Source and licence.} Uniform \v{S}olt\'es' hypergraphs and \v{S}olt\'es' weighted graphs. Version arXiv:2506.07511v1 (CC BY 4.0), distributed under the Creative Commons Attribution 4.0 International licence, as recorded in the arXiv licence metadata for this exact version and shown on the abstract page https://arxiv.org/abs/2506.07511v1. Full rights notice: licenses/arxiv-2506.07511-CC-BY-4.0.txt.\par\smallskip

\noindent\textit{Editorial scope.} Counted unit: the Thm whose source label is \texttt{thm:no4unif\_n9} in the article (source0.tex lines 430--432, source label \texttt{thm:no4unif\_n9}), delivered together with the article's own complete original proof of it (source0.tex lines 434--507, one \texttt{proof} environment of 74 lines). No mathematics is written, repaired or re-derived: statement and proof are the article's own text line for line. Required context retained uncounted: cor:maxWfordiam2 (lines 374--377); lem:Wle44 (lines 383--386); lem:Wbound (lines 410--417). Every definition, display and label of those lines is retained unchanged. Omitted from this package and not redistributed: 1--373, 378--382, 387--409, 418--429, 433--433, 508--648. Nothing inside a retained excerpt is omitted or altered: every retained body is delivered exactly as the article's lines are written, so no omission receipt of any kind accompanies this package. Citations: the retained excerpts carry no citation command at all, so no bibliography entry of the article is transcribed and no reference is redistributed. Reference adaptation: none was needed and none was made; every reference inside the delivered unit resolves inside it, so no unresolved reference or citation is produced. Printed slips kept literal: no printed word, symbol, hypothesis or number of the counted statement or of its proof was altered. Layout and class: the article's own class is \texttt{article} with packages including inputenc, fontenc, wrapfig, amsthm, amsmath, tikz, mathrsfs, amssymb, amsfonts, enumitem, fullpage, caption, hyperref, enumerate; the wrapper is the lane's pinned \texttt{amsart} editorial wrapper at 10pt on A4, which restates the theorem-like environments the retained text needs and adds the article's own macro definitions that the retained text needs (\texttt{\textbackslash abs}, \texttt{\textbackslash diam}), with extra wrapper packages (bm, bbm, dsfont, xspace, cancel, mathtools, cleveref). The delivered numbering is the unit's own. Assets: no retained span contains a figure environment, a graphics-inclusion command, a PDF-page inclusion, a table or a TikZ picture; no image file is copied into this package, the package declares no asset dependency and no image file is redistributed with it. Licence: version arXiv:2506.07511v1 of this article is distributed under the Creative Commons Attribution 4.0 International licence, as recorded in the arXiv licence metadata for this exact version and shown on the abstract page https://arxiv.org/abs/2506.07511v1. A full rights notice is delivered at licenses/arxiv-2506.07511-CC-BY-4.0.txt. Characters: every retained excerpt is pure ASCII, so the delivered engine, XeLaTeX with its default Latin Modern text font, reports no missing character.
\begin{cor}\label{cor:maxWfordiam2}
    A connected $4$-uniform hypergraph $H$ of order $8$ and diameter $2$ satisfies $W(H) \le 41.$
    Equality can only be obtained by hypergraphs of size $3.$
\end{cor}

\begin{lem}\label{lem:Wle44}
    A connected $4$-uniform hypergraph $H$ of order $8$ and diameter $3$ satisfies $W(H) \le 44.$
    Equality is obtained if and only if $H$ has size $3.$
\end{lem}

 \begin{lem}\label{lem:Wbound}
        A connected $4$-uniform hypergraph $H$ of order $8$ and size at least $4$, satisfies either 
        \begin{itemize}
            \item $\diam (H) \le 2$ and $W(H) \le 40$
            \item $\diam (H)  =3$, $W(H)\le 41$ and $H$ has exactly one pair of vertices at distance $3$ 
            \item $\diam(H)  =3$, $40\le W(H) \le 42$ and $H$ has exactly two pairs of vertices at distance $3$ 
        \end{itemize}
    \end{lem}

\begin{thm}\label{thm:no4unif_n9}
     There is no $4$-uniform \v Solt\'es' hypergraph with order $9$.
\end{thm}

\begin{proof}
    Assume the contrary and let $H$ be a $4$-uniform \v Solt\'es' hypergraph with order $9$.

    Observe that $\diam(H)\le 2,$ since $\delta(H)\ge 2$ and no two vertices can have disjoint closed neighbourhoods, since $\abs{N[v]}\ge 5$ for every $v\in V(H).$ 


    \begin{claim}\label{clm:claimdeg2}
        If $\deg(v)=2,$ then $\sigma(v)\ge 10$. 
    \end{claim}
    \begin{claimproof}
        Since $v$ is contained in $2$ edges, it has $r \le 6$ neighbours (equality in the case where $v$ is the only common vertex of these two edges). Thus, $v$ will be at a distance of $1$ from $r$ vertices and $2$ from the remaining $8-r$. This means $\sigma(v)\ge r\times 1+(8-r)\times2=16-r \ge 10$.
    \end{claimproof}

    Let $\deg_H(uv)$ be the number of hyperedges in $E(H)$ containing both $u$ and $v$. If $v$ is fixed, we will call $\deg_H(uv)$ the multiplicity of $u$.

    
    The following claim immediately follows from the fact that at most one $4$-uniform hyperedge can contain four fixed vertices.

\begin{claim}\label{clm:deg(uv)gexuv-1}
    If $\deg_H(uv)\ge 2$, then there is at most one vertex $w \in V \setminus \{u,v\}$ which belongs to all hyperedges containing $u$ and $v$.
\end{claim}

    \begin{claim}\label{clm:claimdeg3}
        If $\deg(v)=3$, then either $\sigma(v)=8$ and the unique $u \in V \setminus v$ with $\deg_H(uv)= 2$ satisfies $d_H(u,v)=d_{H \setminus w}(u,v)$ for every $w \in V \setminus \{u,v\},$ or $\sigma(v)\ge 9$ and there is at least one $u \in V \setminus v$ such that $\deg_H(uv)\ge 2$.
    \end{claim}
    \begin{claimproof}
        If $v$ has eight neighbours (and thus $\sigma(v)=8$), there is exactly one vertex $u$ belonging to two of the hyperedges $e_1, e_2$ containing $v$, and $e_1 \cap e_2 =\{u,v\}$ as all vertices in $V \setminus \{u,v\}$ have exactly one common hyperedge with $v$. This implies $d_H(u,v)=d_{H \setminus w}(u,v)=1$ for every $w \in V \setminus \{u,v\}.$

        If $v$ has less than eight neighbours, then $\sigma(v)\ge 9$ and there is at least one $u \in V \setminus v$ such that $\deg_H(uv)\ge 2$ (by the pigeon hole principle, since the sum of multiplicities of $8$ vertices is $9$).     
    \end{claimproof}


    \begin{claim}\label{clm:claimdeg4}
        If $\deg(v)\ge 4,$ there are at least two neighbours $u_1,u_2 \in V \setminus v$ such that $\deg_H(u_iv)\ge 2$ for $i=1,2$.
    \end{claim}
    \begin{claimproof}
        Assume the contrary, then at most one vertex $u$ of $H \setminus v$ satisfies $\deg_H(uv)\ge 2$. Hence every hyperedge containing $v$ contains at least two vertices $w$ for which $\deg_H(vw)=1.$
        This implies there are at least $8$ vertices sharing only one hyperedge with $v$. Since $\abs{H \setminus v}=8$ and equality is impossible, we conclude.
    \end{claimproof}



Since $H$ is a \v Solt\'es' hypergraph, $H\setminus v$ is a connected $4$-uniform hypergraph of order $8$ and thus satisfies $36=\binom 92 \le W(H)=W(H\setminus v)\le44$ by~\cref{cor:maxWfordiam2} and~\cref{lem:Wle44}.

For every pair of adjacent or non-adjacent vertices $u,v,$ we count the number $x_{uv}$ of vertices $w \in V \setminus\{u,v\}$ for which $d_H(u,v)<d_{H \setminus w}(u,v).$

Observe that by definition
\begin{equation} \label{eq:wiener-index}
    W(H \setminus v_i) = W(H) - \sigma(v_i) + \sum_{\substack{u, w \in V \\ u, w \neq v_i}} \big[d_{H \setminus v_i}(u, w) - d_H(u, w)\big]
\end{equation}
and since $H$ is a \v Solt\'es' hypergraph,
\begin{equation} \label{eq:wiener-index2}
        2W(H)=\sum_{i=1}^{9}\sigma(v_i) = \sum_{i=1}^{9}\sum_{\substack{u, w \in V \\ u, w \neq v_i}} \big[d_{H \setminus v_i}(u, w) - d_H(u, w)\big]
\end{equation}

We now end the proof with a small case-analysis.
If $W(H)=44$, then by~\cref{lem:Wle44}, we know $W(H\setminus v)=44$ implies that $\abs{E(H\setminus v)}=3$ for every $v \in V(H)$ and thus $H$ is $d$-regular, where $d$ satisfies $\frac{9d}{4}-d=3$ and thus $d=\frac{12}{5},$ which is impossible.

If $37\le W(H)\le 43$, we write $W(H)=\binom{9}{2}+r$, where $r=n_2$ satisfies $1\le r\le 7$.

From claims~\ref{clm:claimdeg2},~\ref{clm:deg(uv)gexuv-1},~\ref{clm:claimdeg3} and~\ref{clm:claimdeg4}, we know that $(\sigma(v)-8)+\sum_{u \in N(v)} (2-x_{uv})\ge 2$ for every vertex $v$.
Summing over all $9$ vertices in $V(H)$, results into
$2(36+r)-9\cdot 8+2\sum_{uv \in E(H)}(2-x_{uv})\ge 2\cdot 9$ and thus $\sum_{uv \in E(H)} x_{uv}\le 63-r.$

If $y$ is an upper bound for the number of pairs of vertices at distance $3$ of each other in $H \setminus v_i$, then
$$\sum_{i=1}^{9}\sum_{\substack{u, w \in V \\ u, w \neq v_i}} \big[d_{H \setminus v_i}(u, w) - d_H(u, w)\big] \le 9y+ \sum_{uv \in E(H)} x_{uv}. $$

By~\cref{cor:maxWfordiam2},~\cref{lem:Wle44} and~\cref{lem:Wbound}, we conclude  that~\eqref{eq:wiener-index2} is not satisfied since $2(36+r) > 63-r+18=81-r$ for $4 \le r$ and $2(36+r) > 63-r+9=72-r$ when $1 \le r \le 3$.

For $r=0,$ the last inequality becomes an equality. Every vertex $v$ needs to satisfy $\diam(H \setminus v)=3$ and thus $d_{H \setminus v}(u,w)=3$ for two vertices $u,w \in V.$
Since $W(H)=36$, $\deg(v) \ge 4$ and the equality implies that $u,w$ are the only two vertices sharing multiple hyperedges with $v$ and further that at least two hyperedges contain $\{u,v,w\}$ and thus also $\deg_H(vw) \ge 2.$
Since the properties for $u,w,v$ are the same, there should be three disjoint hyperedges for which $\abs{ e \cap \{u,v,w\}}=1$, which is impossible. We conclude that no equality is possible in this case.
% The $9$ vertices can be grouped in three triplets $\{u_i,v_i,w_i\}$, $i \in \{1,2,3\}$, such that there are precisely three hyperedges $e$ satisfying $\abs{ e \cap \{u_i,v_i,w_i\}}=1$.
\end{proof}

\end{document}
